\documentclass[12pt]{article}
\usepackage[utf8]{inputenc}
\usepackage{amsmath}
\usepackage{graphicx}
\usepackage{hyperref}
\usepackage{geometry}
\geometry{a4paper, margin=1in}

\title{Omission halts local field potential structure and reshapes spiking across the cortical hierarchy}
\author{Hamed Nejat \\ Dept. of Psychology, Vanderbilt University \\ Advisor: Dr. Andre Bastos, Vanderbilt Brain Institute}
\date{\today}

\begin{document}

\maketitle

\section*{Abstract}
Predictive coding proposes that the brain continuously compares expected input with actual input and tries to minimize the mismatch between them. In neurophysiology, the predictive routing framework proposes that this process is expressed through oscillatory dynamics across the cortical hierarchy (V1, V2, V3d, V3a, V4, MT, MST, TEO, FST, FEF, PFC)...

\section{Introduction}
[Insert full original Introduction here]

\section{Methods}
Neural data were acquired from multi-area, dense laminar macaque electrophysiology experiments (V1, V2, V3d, V3a, V4, MT, MST, TEO, FST, FEF, and PFC). 
[Insert expanded technical methods from \texttt{D:\textbackslash drive\textbackslash omission\textbackslash docs\textbackslash methods\textbackslash}]

\section{Results}
Omission-sensitive units are a selective minority across the 11-area cortical hierarchy.
[Insert expanded results from \texttt{D:\textbackslash drive\textbackslash omission\textbackslash docs\textbackslash results\textbackslash}]

\section{Discussion}
[Insert original Discussion text]

\section{References}
[Insert consolidated Bibliography]

\end{document}
